\section{Results}
SessionKV-DR provides the core empirical evidence. Its design links success directly to session memory operations, making it suitable for mechanism validation.

Across multi-seed comparisons, full ALC improves delayed session recall relative to baseline controls in the reported slices. Competitive no-write cases are kept to show where write decisions and interference remain limiting factors.

TinyQA serves as a secondary check and is interpreted conservatively.

\subsection{Primary evidence: SessionKV-DR}
\begin{table}[h]
\centering
\begin{tabular}{llll}
\toprule
Delay/Facts & Baseline & ALC no-write & ALC full \\
\midrule
d=0, f=4 & 0.125 ± 0.000 (95% CI ± 0.000, n=5) & 0.125 ± 0.000 (95% CI ± 0.000, n=5) & 1.000 ± 0.000 (95% CI ± 0.000, n=5) \\
d=2, f=4 & 0.125 ± 0.000 (95% CI ± 0.000, n=5) & 0.125 ± 0.000 (95% CI ± 0.000, n=5) & 0.875 ± 0.000 (95% CI ± 0.000, n=5) \\
d=6, f=4 & 0.125 ± 0.000 (95% CI ± 0.000, n=5) & 0.125 ± 0.000 (95% CI ± 0.000, n=5) & 0.719 ± 0.000 (95% CI ± 0.000, n=5) \\
d=12, f=4 & 0.125 ± 0.000 (95% CI ± 0.000, n=5) & 0.125 ± 0.000 (95% CI ± 0.000, n=5) & 0.469 ± 0.000 (95% CI ± 0.000, n=5) \\

\bottomrule
\end{tabular}
\caption{SessionKV-DR delayed recall (mean $\pm$ std, 95\% CI, n=5).}
\end{table}

\begin{table}[h]
\centering
\begin{tabular}{llll}
\toprule
Metric & Baseline & ALC no-write & ALC full \\
\midrule
Overwrite accuracy & 0.125 ± 0.000 (95% CI ± 0.000, n=5) & 0.125 ± 0.000 (95% CI ± 0.000, n=5) & 0.375 ± 0.000 (95% CI ± 0.000, n=5) \\
Stale confusion rate & 0.125 ± 0.000 (95% CI ± 0.000, n=5) & 0.125 ± 0.000 (95% CI ± 0.000, n=5) & 0.625 ± 0.000 (95% CI ± 0.000, n=5) \\
Persistence recall acc & 0.000 ± 0.000 (95% CI ± 0.000, n=5) & 0.000 ± 0.000 (95% CI ± 0.000, n=5) & 1.000 ± 0.000 (95% CI ± 0.000, n=5) \\

\bottomrule
\end{tabular}
\caption{SessionKV-DR overwrite and persistence metrics (mean $\pm$ std, 95\% CI, n=5 where available).}
\end{table}

\begin{table}[h]
\centering
\begin{tabular}{llll}
\toprule
Setting & Baseline & ALC no-write & ALC full \\
\midrule
facts=1, delay=6 & 0.125 ± 0.000 (95% CI ± 0.000, n=5) & 0.125 ± 0.000 (95% CI ± 0.000, n=5) & 1.000 ± 0.000 (95% CI ± 0.000, n=5) \\
facts=4, delay=6 & 0.125 ± 0.000 (95% CI ± 0.000, n=5) & 0.125 ± 0.000 (95% CI ± 0.000, n=5) & 0.719 ± 0.000 (95% CI ± 0.000, n=5) \\
facts=8, delay=6 & 0.125 ± 0.000 (95% CI ± 0.000, n=5) & 0.125 ± 0.000 (95% CI ± 0.000, n=5) & 0.719 ± 0.000 (95% CI ± 0.000, n=5) \\
distractors=0, facts=4 & 0.125 ± 0.000 (95% CI ± 0.000, n=5) & 0.125 ± 0.000 (95% CI ± 0.000, n=5) & 1.000 ± 0.000 (95% CI ± 0.000, n=5) \\
distractors=4, facts=4 & 0.125 ± 0.000 (95% CI ± 0.000, n=5) & 0.125 ± 0.000 (95% CI ± 0.000, n=5) & 0.844 ± 0.000 (95% CI ± 0.000, n=5) \\
distractors=12, facts=4 & 0.125 ± 0.000 (95% CI ± 0.000, n=5) & 0.125 ± 0.000 (95% CI ± 0.000, n=5) & 0.344 ± 0.000 (95% CI ± 0.000, n=5) \\
distractors=24, facts=4 & 0.125 ± 0.000 (95% CI ± 0.000, n=5) & 0.125 ± 0.000 (95% CI ± 0.000, n=5) & 0.125 ± 0.000 (95% CI ± 0.000, n=5) \\

\bottomrule
\end{tabular}
\caption{SessionKV-DR capacity/interference sweep.}
\end{table}

\subsection{Supporting evidence: TinyQA}
\begin{table}[h]
\centering
\begin{tabular}{llll}
\toprule
Delay & Baseline & ALC no-write & ALC full \\
\midrule
0 & 0.033 ± 0.000 (95% CI ± 0.000, n=5) & 0.033 ± 0.000 (95% CI ± 0.000, n=5) & 0.033 ± 0.000 (95% CI ± 0.000, n=5) \\
4 & 0.033 ± 0.000 (95% CI ± 0.000, n=5) & 0.033 ± 0.000 (95% CI ± 0.000, n=5) & 0.067 ± 0.000 (95% CI ± 0.000, n=5) \\
8 & 0.033 ± 0.000 (95% CI ± 0.000, n=5) & 0.033 ± 0.000 (95% CI ± 0.000, n=5) & 0.050 ± 0.000 (95% CI ± 0.000, n=5) \\

\bottomrule
\end{tabular}
\caption{TinyQA delayed QA accuracy (supporting evidence only).}
\end{table}

